\section{Concrete tensor identities on the square}\label{sec:identity-catalogue}
We now evaluate several identities with their selector costs included. Throughout this section the square has side length one, its vertices are $1,2,3,4$ in cyclic order, and its centre is $0$. Thus the diagonal length is $\sqrt2$ and $e'(0,i)=1/\sqrt2$. By Corollary~\ref{cor:price-attained}, a price representation can be chosen as $e=c\,d_\psi$ or $e'=c\,d_\psi$ with $\psi$ an ordinary normalized spectral point. Fix this same $\psi$ throughout each test, and write
\[
 \gamma(P)=c\log_2\psi(U_2(P)).
\]
The quantity $\gamma(P)$ is the scaled cost of a shared selector. Every pairwise distance within $P$ is at most $\gamma(P)$, and every tree joining $P$, with available hubs, gives an upper bound for $\gamma(P)$. A pure rank bound for a graph tensor therefore supplies an inequality involving $\gamma(P)$; it does not by itself specify this quantity.

\subsection{Finite instances of the crossing identity}
The determinant identity in \cite[Proposition B.2, pp.~23--24]{Anthropic2026} is indexed by Frobenius symbols of fixed size and energy. We evaluated some small instances, keeping both adjacent-pair flattenings of every coefficient tensor. The terminal positions in this subsection are the unit-circle positions of Proposition~\ref{prop:optimal-crossing-hubs}; their distances are $\sqrt2$ times those of the side-one square.

For each coefficient tensor $T_j$, let $r_j$ be the smaller of its ranks across the two adjacent-pair cuts, and put $r_* =\max_jr_j$. A rank factorization realizes $T_j$ on the corresponding two-hub tree with bridge dimension $r_j$. After taking tensor powers, integral cable widths incur an arbitrarily small relative overhead. With $q=2^m$, the best worst-case cost of these certificates, divided by $m$, is
\[
 F\!\left(\frac{\log_2r_*}{\log_2q}\right).
\]
The selector charge divided by $m$ is $4\log_2 A/\log_2 q$, where $A$ is the number of nonzero coefficient tensors. This uses the full number of terms, not the number of distinct numerical ranks.

\begin{center}
\begin{tabular}{@{}ccccrrcc@{}}
\toprule
Field & $q$ & Size & Energy & $A$ & $r_*$ & $\log_qr_*$ & Cost$/m$\\
\midrule
$\mathbb Q$ & $2$ & $1$ & $4$ & $9$ & $4$ & $2$ & $4$\\
$\mathbb F_2$ & $2$ & $1$ & $4$ & $8$ & $2$ & $1$ & $\sqrt2+\sqrt6$\\
$\mathbb F_2$ & $4$ & $2$ & $12$ & $25$ & $11$ & $1.729716$ & $4$\\
$\mathbb F_2$ & $8$ & $2$ & $14$ & $29$ & $54$ & $1.918296$ & $4$\\
$\mathbb F_2$ & $16$ & $2$ & $18$ & $37$ & $253$ & $1.995748$ & $4$\\
\bottomrule
\end{tabular}
\end{center}

The ranks and term counts in the table are exact integers; the logarithms are rounded to six decimal places. For $q=4$ and $q=16$ the chosen symbols form specified subsets of classes of sizes $5$ and $20$, respectively; the $q=8$ case uses its entire class of size $8$. These choices are recorded with the coefficient-rank profiles. The rational calculation expands the determinant exactly and checks its reflection identity. The characteristic-two calculation uses interpolation over an explicit finite extension, checks that every coefficient belongs to $\mathbb F_2$, checks the sum identity against the independently formed crossing tensor, and then uses binary row reduction. The size-one characteristic-two result agrees with reduction of the rational calculation.

None of these cases gives a strict shortening with these certificates. For $q=4,8,16$ even the alternative networks have worst-case cost $4m$, before the positive selector cost is added. For $q=2$ in characteristic two, the alternative cost is approximately $3.863704m$, but the selector charge is $12m$. The conclusions concern the displayed coefficient tensors and their two-hub rank factorizations. They do not bound different layouts, larger symbol classes, different decompositions, or the asymptotic identity in general.

\subsection{Two and three crossing wires, and the four sides}
For the two-term Pl\"ucker identity from Section~\ref{sec:two-hub-cost}, the side-one square gives
\begin{equation}\label{eq:plucker-price-test}
 2\sqrt2m\leq 2m+m\gamma(\{1,2,3,4\}).
\end{equation}
Indeed the $m$th power has $2^m$ alternatives, each a product of adjacent wires of cost $2m$. Routing the selector on a corner spanning tree bounds the right side by $5m$. With the centre available it is bounded by $(2+2\sqrt2)m$. Both tests are satisfied. The same identity on either pair among three crossing wires, tensored with the third unchanged wire, has the same cost difference: the unchanged wire cancels from both sides. This observation addresses these tensor-product extensions only; an identity coupling all three crossings can behave differently.

One can also couple all three crossings. Write $[ij]=J_{ij}$ for the alternating dimension-two wire. Repeated use of the same Pl\"ucker relation gives the integral identity
\begin{align}\label{eq:three-crossings}
 [14][25][36]={}&[12][36][45]+[12][34][56]+[14][23][56]\nonumber\\
 &+[16][23][45]+[16][25][34].
\end{align}
For a direct verification, first expand $[14][25]$ as $[12][45]+[15][24]$, then expand $[15][36]$ as $[13][56]+[16][35]$. Finally expand $[13][24]$ and $[24][35]$. These four applications produce exactly the five terms displayed, over every field. On the unit regular hexagon, with vertices in cyclic order, the original three diameters have cost $6$ and the five alternatives have respective costs $4,3,4,3,4$. Their centre selector costs $6\log_2 5$, so this does not shorten the original wires. For the orthogonal terminals $e_1,e_2,e_3,-e_1,-e_2,-e_3$ of $C_{3,1}$, the five costs instead alternate between $2+2\sqrt2$ and $3\sqrt2$ in the same pattern; the worst alternative plus its selector again exceeds $6$. The $m$th tensor power multiplies every cost and the logarithm of the number of alternatives by $m$, and therefore gives no improvement to either test.

There is also a direct four-side test. Apply the Pl\"ucker relation, after relabelling, to each of the two pairs of opposite sides. Their product expresses the four-side tensor as four alternatives. In one alternative both pairs become diagonals, so the worst cost is $4\sqrt2m$, while the original cost is $4m$. The selector has $4^m$ labels. Thus this construction already loses before the selector charge $2m\gamma(\{1,2,3,4\})$ is included.

The ordinary rank and border rank of the dimension-$q$ four-side tensor are both $q^4$. To see this, flatten the vertices $\{1,3\}$ against $\{2,4\}$: every side crosses the cut, giving matrix rank $q^4$. Its defining expansion has exactly $q^4$ simple terms. A rank-one cover with this count gives
\[
 4m\leq4m\gamma(\{1,2,3,4\}),
\]
which is weaker than the lower bound on $\gamma$ supplied by any diagonal. Proposition~\ref{prop:single-layout} applies much more generally to every single restriction or degeneration between a wire graph and one network: such a realization cannot reduce Euclidean cost, regardless of how many of its wires cross in a drawing. A sum requires its own quantitative analysis.

\subsection{The W tensor and matrix multiplication}
On three parties put
\[
 W=e_1\otimes e_0\otimes e_0+
 e_0\otimes e_1\otimes e_0+
 e_0\otimes e_0\otimes e_1.
\]
The coefficient of $t$ in $(e_0+te_1)^{\otimes3}-e_0^{\otimes3}$ is $W$. This gives a degeneration from the dimension-two unit tensor to $W$ in every characteristic. Contracting any one party of $W$ with $e_0^*$ gives a rank-two wire on the other two parties. Therefore, for any triple $P$,
\begin{equation}\label{eq:w-price-test}
 \max_{i,j\in P}d(i,j)
 \leq c\log_2\psi(W_P)
 \leq\gamma(P)
 \leq\min_{\tau\text{ joining }P}\sum_{ij\in E(\tau)}d(i,j).
\end{equation}
For three consecutive corners the outer interval is $[\sqrt2,2]$. For the centre and two opposite corners it collapses to the single value $\sqrt2$. Both are consistent. In particular the latter equality gives information about a mixed tensor, but it is not a contradiction.

The triangle graph tensor is the matrix multiplication tensor. If an asymptotic rank estimate of exponent $\omega$ is used on three parties at locations $x,y,z$, its price inequality is
\begin{equation}\label{eq:mm-price-test}
 d(x,y)+d(y,z)+d(z,x)\leq\omega\gamma(\{x,y,z\}).
\end{equation}
A flattening forces $\omega\geq2$. Every tree joining three points has length at least half their perimeter: each pairwise path has length at least the corresponding distance, and the three paths together traverse each edge of the minimal connecting subtree twice. Consequently, supplying the unit tensor through such a tree and then applying any rank algorithm with exponent at least two cannot shorten the triangle wires. For the square triples this is also visible directly: the perimeter is $2+\sqrt2$ for three corners, and $2\sqrt2$ for the centre and opposite corners. Thus even the newly claimed exponent $9/4$, whose untwisting mechanism is discussed in \cite[Sections 2--3]{Li2026}, does not yield a contradiction through this triangle-and-tree test.

\subsection{Four-party graph bounds and a remaining interval}
The subrank calculation needed here has a short proof specialized to the cycle. It is a special case of the graph distillation theorem of Vrana and Christandl \cite[Theorem 1, p.~3, and proof in Section 2]{VC2017}.

\begin{lemma}[Subrank of the four-side tensor]\label{lem:cycle4-subrank}
Over every field, the asymptotic subrank of the dimension-two four-cycle tensor is $4$.
\end{lemma}
\begin{proof}
Give the four cyclic edges the integer vectors
\[
 c_1=(1,0),\quad c_2=(1,1),\quad
 c_3=(0,1),\quad c_4=(1,-1).
\]
Opposite, nonincident edges have orthogonal vectors, and any two of the four vectors are linearly independent over $\mathbb Q$. For edge indices $0\leq i_j<q$, the vector $\sum_jc_ji_j$ takes fewer than $9q^2$ possible values. Some value $g$ therefore has at least $q^2/9$ preimages. Multiply the term indexed by $(i_1,i_2,i_3,i_4)$ by the formal power
\[
 t^{\|\sum_jc_ji_j-g\|^2}.
\]
This is a product of local diagonal changes: in the expanded square, every mixed term involving nonincident edges vanishes, and every other term depends only on indices visible together at one vertex. Linear, square, and constant terms can also be assigned to vertices. Local exponents may be negative, but their sum is a nonnegative integer, so the constant coefficient is a valid degeneration. The exponents are integers; their role is unchanged by the characteristic of the ground field.

The surviving tensor is a unit tensor. Indeed the two indices visible at any vertex uniquely determine the other two through the equation $\sum_jc_ji_j=g$, because the two remaining coefficient vectors are independent. No local basis vector therefore occurs in two surviving terms.

Fix $q$ first, and write $s\geq q^2/9$ for the size of this unit tensor.
The $N$th tensor power of the degeneration has local polynomial degree
$O_q(N)$ after shifting the Laurent exponents. Extracting its designated
coefficient expresses the unit tensor of size $s^N$ as a sum of
$O_q(N^3)$ restrictions of the $N$th power of the cycle tensor: choose
the first three local degrees, after which the fourth is determined.
This coefficient argument uses no division and works over finite fields.
A spanning tree of one dimension-$q$ cycle supplies a unit tensor of
size $q$. Thus $O_q(\log N)$ extra cycle copies supply the polynomial
selector, giving an actual restriction from $N+O_q(\log N)$ copies
to the unit tensor of size $s^N$. Taking $N$th roots with $q$ fixed
proves asymptotic subrank at least $q^2/9$ for that cycle.

Only now put $q=2^m$. This cycle is the $m$th tensor power of the
dimension-two cycle, so its asymptotic subrank gives the lower bound
$(4^m/9)^{1/m}$ for the latter. Letting $m$ tend to infinity gives
$4$. The opposite upper bound is the single-vertex flattening, of
rank $4$ for the dimension-two cycle.
\end{proof}

Suppose a normalized four-party spectral point has wire metric $d_\psi=\lambda e$. Lemma~\ref{lem:cycle4-subrank} gives $4\lambda\geq2$. The complete-graph asymptotic-rank bound in the recent preprint \cite[Theorem 1.1]{Zhou2026}, taken as an external input, gives
$(4+2\sqrt2)\lambda\leq4$. Together they leave the nonempty interval
\begin{equation}\label{eq:square-scale-interval}
 \frac12\leq\lambda\leq2-\sqrt2.
\end{equation}
Equivalently, for a scaled spectrum on a larger party set whose four corner distances are $e$, these estimates say
\[
 1+\frac1{\sqrt2}\leq\gamma(\{1,2,3,4\})\leq2.
\]
The subrank inequality also sharpens the Pl\"ucker test \eqref{eq:plucker-price-test} to the upper estimate $2\sqrt2m\leq4m$, which remains satisfied.

All linear inequalities on normalized wire exponents valid throughout the spectrum are consistent with $\lambda=1/2$. Here is an explicit certificate. Let $\delta_S$ denote the cut metric of a nonempty proper subset $S$ of the four corners. Then
\begin{equation}\label{eq:square-convex-cuts}
 \frac e2=
 \frac{2-\sqrt2}{4}\sum_{i=1}^4\delta_{\{i\}}
 +\frac{\sqrt2-1}{2}
   \bigl(\delta_{\{1,2\}}+\delta_{\{2,3\}}\bigr).
\end{equation}
The coefficients are nonnegative and sum to one. Each cut metric is the wire metric of the normalized spectral point given by the corresponding flattening rank. To check the identity, a side is separated by two singleton cuts and one of the last two cuts; a diagonal is separated by two singleton cuts and both of the last two cuts. Thus $e/2$ lies in the convex hull of normalized spectral wire metrics. This does not assert that it is itself a spectral wire metric: convex combinations of those metrics need not be realized by a spectral point. It does show that graph rank and subrank bounds, each of which is a linear inequality in the wire exponents, cannot settle the square question by themselves.

The same normalized candidate survives when the centre is adjoined:
\begin{equation}\label{eq:centred-square-convex-cuts}
 \frac{e'}2=
 \frac{2-\sqrt2}{4}\sum_{i=1}^4\delta_{\{i\}}
 +\frac{\sqrt2-1}{4}
  \bigl(\delta_{\{1,2\}}+\delta_{\{2,3\}}+
        \delta_{\{3,4\}}+\delta_{\{4,1\}}\bigr).
\end{equation}
These are cuts on all five parties, with the centre outside each displayed set. Their coefficients again sum to one. If the two coefficients are denoted by $s,t$, respectively, a radial pair has value $s+2t=1/(2\sqrt2)$, a side has value $2s+2t=1/2$, and a diagonal has value $2s+4t=1/\sqrt2$. Thus every affine linear inequality valid on normalized spectral wire metrics also passes this five-party candidate.

There is a sharper upper bound on the normalization scale from graph subrank alone. We state exactly what that bound optimizes.
\begin{proposition}[The strongest graph-subrank scale bound on the five points]\label{prop:centred-subrank-scale}
Assume the graph minimum-cut subrank formula of \cite[Theorem 1 and Section 3]{VC2017}. If $e'=c\,d_\psi$ for a normalized spectral point on the five parties, then
\[
 c\leq\frac{3+\sqrt2}{2},\qquad
 d_\psi=\lambda e'\ \Longrightarrow\
 \lambda\geq\frac{2}{3+\sqrt2}.
\]
The upper bound on $c$ is optimal among bounds obtained by dividing a nonnegative weighted graph's Euclidean wire cost on these five points by its minimum cut weight.
\end{proposition}
\begin{proof}
For any graph with integral edge multiplicities $m_{ij}$, the cited subrank formula and spectral monotonicity give
\[
 c\leq\frac{\sum_{i<j}m_{ij}e'(i,j)}
                   {\operatorname{MinCut}(m)}.
\]
Rational approximation extends this optimization to nonnegative real weights with positive minimum cut. Take multiplicity three on each side and two on each radial edge. Its minimum cut is eight and its length is $12+4\sqrt2$, giving the stated bound.

For optimality normalize the minimum cut to at least one and average the weights over the dihedral symmetries of the square. Every cut inequality is preserved by this averaging, and the cost is unchanged. Let $a,b,r$ be the resulting weights on each side, each diagonal, and each radial edge. The single-corner and centre cuts give
\[
 2a+b+r\geq1,\qquad4r\geq1.
\]
Consequently the cost is bounded below by
\begin{align*}
 4a+2\sqrt2b+2\sqrt2r
 &=2(2a+b+r)+(2\sqrt2-2)(b+r)\\
 &\geq2+\frac{2\sqrt2-2}{4}
   =\frac{3+\sqrt2}{2}.
\end{align*}
The choice $a=3/8$, $b=0$, $r=1/4$ attains this value. To verify all cut constraints, replace a cut by its complement so that it excludes the centre. A single corner, two adjacent corners, two opposite corners, three corners, and all four corners have respective cut weights $1,5/4,2,3/2,1$. Hence the candidate is feasible and optimal. This proves optimality of this family of bounds, not existence of a spectrum attaining the bound.
\end{proof}

\subsection{A shared-mode splitting obstruction and its precise limitation}
The untwisting lemma of \cite[Theorem 2.1, pp.~5--8]{Li2026} concerns a semidirect sum with one shared tensor leg. Its proof uses the structure tensor of $\kk[t]/(t^p)$, whose border rank is $p$. Copying that shared output coordinate into two parties changes this cost substantially.

\begin{proposition}[The cost of a copied polynomial output]\label{prop:copied-output}
For any field $\kk$ and integer $p\geq1$, set
\[
 T_p=\sum_{\substack{j,k\geq0\\j+k<p}}
 e_j^{(a)}\otimes e_k^{(b)}\otimes
 e_{j+k}^{(c)}\otimes e_{j+k}^{(d)}.
\]
Its tensor rank over $\kk$ and its border rank over an algebraic closure of $\kk$ are both $p(p+1)/2$.
\end{proposition}
\begin{proof}
Flatten the parties $a,c$ against $b,d$. A nonzero matrix entry has a common output index $\ell=j+k$ on $c,d$. After grouping rows and columns by this index, the nonzero blocks are indexed by $0\leq\ell<p$. The block for $\ell$ has rows $j=0,\ldots,\ell$, columns $k=0,\ldots,\ell$, and entry one precisely when $j+k=\ell$. It is a reversal permutation matrix of size $\ell+1$. These disjoint blocks have total matrix rank $\sum_{\ell=0}^{p-1}(\ell+1)=p(p+1)/2$ over every field. Matrix rank gives both the rank lower bound and the border-rank lower bound. The displayed tensor expression has exactly that many simple terms and gives the matching upper bound.
\end{proof}

The proposition rules out only the literal copied-output auxiliary tensor. It does not rule out simultaneous untwisting with several shared legs. In fact a different construction in \cite[Theorem 3.4 and Corollary 3.5]{Zhou2026} achieves a linear auxiliary cost. Because this distinction matters for possible four-party identities, we give the construction explicitly.

Let $A=A_0\boxplus_e\cdots\boxplus_e A_{M-1}$ be a $k$-party partial direct sum of nonzero blocks, where $e=(k-1,k)$: the last two parties have disjoint sectors labeled by $h$, and the first $k-2$ parties share their coordinate spaces. Put $D=1+(k-1)(M-1)$. Zhou's construction gives the formal degeneration
\begin{equation}\label{eq:several-shared-legs}
 U_D([k])A\ \unrhd\
 \bigoplus_{h=0}^{M-1}\bigl(A_h\langle M\rangle_e\bigr).
\end{equation}
Here $\unrhd$ means that the tensor on the left degenerates to the one on the right.

To prove it, use the $k$-party polynomial multiplication tensor
\[
 C=\sum_{0\leq g_1,\ldots,g_{k-2},u<M}
 e_{g_1}\otimes\cdots\otimes e_{g_{k-2}}\otimes e_u
 \otimes e_{g_1+\cdots+g_{k-2}+u}.
\]
It degenerates from $U_D([k])$: evaluate the input polynomials at $D$ distinct points over the rational function field $\kk(z)$, and interpolate their product, of degree at most $D-1$. This is a rank-$D$ expression over $\kk(z)$. Clear its denominators and take the lowest nonzero power of $z$ to obtain a formal degeneration over $\kk$ itself.

In $A\otimes C$, assign weight $g_j^2$ to shared party $j$, weight $2hu+(k-2)h^2$ to party $k-1$, and weight $-2hv$ to party $k$, where $v=\sum_jg_j+u$. The total weight on a supported entry is
\[
 \sum_{j=1}^{k-2}(g_j-h)^2.
\]
It is nonnegative and vanishes precisely when all the tentative shared labels equal $h$. In this surviving part $u$ is unrestricted and $v=(k-2)h+u$. The label $h$, now visible at every party, separates the summands completely; relabelling $v$ locally retains an $M$-dimensional wire on $e$. This proves \eqref{eq:several-shared-legs}. All squares are integer exponents of a formal variable, so no characteristic restriction is involved.

For a normalized spectral point, put $p=d_\psi(e)$. Apply \eqref{eq:several-shared-legs} to the blocks of a fixed multiplicity type in $A^{\otimes N}$. There are $M_N=N!/\prod_h(N\theta_h)!$ such blocks, each of spectral value $\prod_h\psi(A_h)^{N\theta_h}$. Since the auxiliary dimension is at most $(k-1)M_N$, additivity and monotonicity give
\[
 (k-1)\psi(A)^N\geq
 M_N^{p}\prod_h\psi(A_h)^{N\theta_h}.
\]
Take $N$th roots, use Stirling's formula, and approximate arbitrary probability vectors by rational ones. In base-two notation this proves
\begin{equation}\label{eq:shared-entropy}
 \log_2\psi(A)\geq pH_2(\theta)+
                  \sum_h\theta_h\log_2\psi(A_h).
\end{equation}
This is Zhou's entropy inequality. It is a lower bound for a partial direct sum, whereas Lemma~\ref{lem:terminal-selector} bounds an ordinary sum from above. Using both on a suitable mixed tensor is a possible route beyond the linear graph inequalities. An ordinary sum identity does not automatically have the required disjoint sectors on two specified parties; that structural condition must be checked for any proposed application.

\subsection{Applying border-rank and pencil methods}\label{sec:border-pencil-tests}
The border-rank and residual-subspace methods provide three different kinds of information: a lower-bound map, a reduction to a smaller residual tensor, and an explicit border degeneration. Their geometric uses require keeping these roles distinct.

For the multiple exterior map in \cite[sections ``The multiple exterior map for the bound fifty-two'' and ``The exterior map used for the uniform bound'']{SuppliedBorder2026}, wedge multiplication by a nonzero vector in dimension $a$ has rank $\binom{a-1}{p}$ in exterior degree $p$. Thus two such wedge maps have pure four-tensor cost $\binom{a-1}{p}\binom{b-1}{r}$. The seven-dimensional degree-three example has cost $20^2=400$. Linearity and closed determinantal loci give the corresponding border-rank inequality. That cost concerns four separate factors; grouping two of them changes what counts as a pure summand and does not preserve this calibration.

For the two-hub layouts actually tested in this paper, ordinary matrix rank already gives an exact criterion.
\begin{lemma}[The bridge dimension for two hubs]\label{lem:twohub-rank}
Let a four-party tensor have terminal dimensions at most $Q$. A layout with dimension-$Q$ access wires, one hub adjacent to $a,b$, one adjacent to $c,d$, and a bridge of dimension $R$ realizes the tensor if and only if its $ab\mid cd$ flattening has rank at most $R$. Allowing a border realization on this fixed layout does not change the criterion.
\end{lemma}
\begin{proof}
Contracting each hub with its incident local tensors factors the flattening through the bridge space, proving necessity. Conversely a rank factorization
$T=\sum_{i=1}^R A_i^{ab}\otimes B_i^{cd}$
specifies the two hub tensors; the access wires transmit the terminal coordinates. The rank-at-most-$R$ matrix locus is closed, so a degeneration cannot enlarge it.
\end{proof}

The multidimensional substitution inequality in \cite[section ``The all-field multidimensional substitution statement'']{SuppliedPencils2026} keeps all residual factors separate. If the first flattening image is $U$ of dimension $a$, its statement is
\[
 \underline{\operatorname R}(T)\geq
 a-j+\min_{L\in\operatorname{Gr}(j,U)}
       \underline{\operatorname R}(P_L).
\]
The proof uses a closed flag incidence and its proper projection, so it remains valid over an algebraically closed field in positive characteristic. The coordinate reductions and adjacent-pair calculations in the section ``Exact two-dimensional residual minima'' give exact two-dimensional residual minima $3$ for the row determinant/permanent families and $6$ for the ordered polynomial coefficient families. For $n=4$, substitution therefore yields $7$ and $40$, respectively. These are lower bounds for those coefficient tensors, not wire identities on the square.

Exact integer pencil certificates give the following counts. Among the $15$ unordered row-edge pairs on four labels, $12$ have support rank three and $3$ rank four. Among the $630$ unordered pairs of coordinate quadrics, $144$ have support rank six, $288$ rank seven, and $198$ rank eight. For both determinant and permanent signs, selecting one matching column for each support variable gives a diagonal signed unit minor. The two-dimensional Koszul map is exactly the concatenated slice matrix, up to a column permutation and signs, so no unproved generic-rank assertion enters these counts.

The following lemma gives an exact obstruction to copying the shared selector leg of one of the border constructions.
\begin{lemma}[Copying the slice label]\label{lem:copied-slices}
Let $M_1,\ldots,M_h\in B\otimes C$ be matrices, and let $(a_i)$ and $(d_i)$ be independent labels. Over every field,
\[
 T=\sum_{i=1}^h a_i\otimes M_i\otimes d_i
 \quad\Longrightarrow\quad
 \operatorname R(T)=\underline{\operatorname R}(T)
 =\sum_i\operatorname{rank}(M_i),
\]
where border rank is taken geometrically over the algebraic closure.
\end{lemma}
\begin{proof}
The flattening $AB\mid CD$ is the direct sum of the matrices $M_i$, so its rank is the stated sum. This is a lower bound for rank and border rank over every extension field. Decomposing each $M_i$ into its matrix-rank number of simple tensors gives exactly that many pure four-tensors, proving the matching upper bound.
\end{proof}

For example, the star tensor is
\[
 S_{k,\delta}=\sum_{i=1}^k a_i\otimes
 (e_0\otimes e_i+\delta e_i\otimes e_0),\qquad\delta=\pm1.
\]
Its flattening at the second party has rank $k+1$. The exact curve
\[
 t^{-1}\left[\sum_i a_i\otimes(e_0+\delta t e_i)
                  \otimes(e_0+t e_i)
 -(\sum_i a_i)\otimes e_0\otimes e_0\right]
 =S_{k,\delta}+\delta t\sum_i a_i\otimes e_i\otimes e_i
\]
has $k+1$ pure terms and proves its border rank is $k+1$ in every characteristic. Copying the $a_i$ label to a fourth party gives rank $2k$ by Lemma~\ref{lem:copied-slices}. Likewise, each slice of the anchored tensor is a nonsingular matching on four independent variables. Its copied version has rank $4h$, whereas the original three-party border rank is $h+r+s+1$ \cite[``anchored collections have explicit exact border rank'']{SuppliedPencils2026}. The exceptional permanent finite-difference constructions have five or six such slices and upper border bounds $11$ or $12$ \cite[``finite differences for the exceptional permanent subspaces'']{SuppliedPencils2026}; after copying, their ranks are $20$ or $24$.

These star degenerations give explicit valid mixed-tensor price inequalities. Put $v_k=c\log_2\psi(S_{k,\delta})$ on parties $a,b,c$, and use the same $\gamma$ as in \eqref{eq:w-price-test}. Contracting the $e_0$ coordinate at either last party leaves a dimension-$k$ wire from the first party to the other. Selecting a first-party label leaves a rank-two wire between the last parties. The border curve gives the upper bound, by coefficient extraction on tensor powers and its polynomial overhead. Hence
\begin{equation}\label{eq:star-price-test}
 \max\{(\log_2 k)d(a,b),(\log_2 k)d(a,c),d(b,c)\}
 \leq v_k\leq(\log_2(k+1))\gamma(\{a,b,c\}).
\end{equation}
Supplying the selector on a tree bounds the right endpoint by its tree length times $\log_2(k+1)$. For three consecutive square corners, with the first party at the middle corner, the resulting interval is
$[\max\{\log_2 k,\sqrt2\},\,2\log_2(k+1)]$.
For the centre as first party and two opposite corners it is
$[\max\{(\log_2 k)/\sqrt2,\sqrt2\},\,\sqrt2\log_2(k+1)]$.
Both intervals are nonempty for every $k\geq1$. Thus these particular bounds do not exclude $e$ or $e'$.

\paragraph{The exceptional five-plane, checked in characteristic two.}
The necessary-prolongation calculation for the exceptional permanent plane \cite[section ``A characteristic-free obstruction for the exceptional permanent plane'']{SuppliedPencils2026} admits the following characteristic-two certificate. Let $A$ have basis $c_0,\ldots,c_4$ and let $B,C$ have variables $a,b,c,d,e,f,g,h,i$. The five symmetric ordered slices are
\[
 ae+bd,\quad af+cd,\quad bf+ce,\quad ah+bg,\quad ai+cg,
\]
where each product $xy$ denotes $x\otimes y+y\otimes x$. Their tensor is denoted $T_5$. Let $U=T_5(C^*)\subset A\otimes B$, and let $Z_s(X)=c_s\wedge X$. Define
\[
 K=[Z_0U\mid\cdots\mid Z_4U],\qquad
 E_0=[U\mid c_2\otimes a].
\]
Over $\mathbb F_2$, choose $L$ as a row-reduced basis of the left kernel of $K$, put $H$ equal to the vertical stack of $LZ_s$, and put $R=\Phi_B(\operatorname{span}E_0)$, where
$\Phi_B:B\otimes(A\otimes B)\to A\otimes\Lambda^2B$ wedges in the two $B$ slots. Exact elimination gives
\[
 \operatorname{rank}(U,K,L,H,E_0,R)=(9,40,50,35,10,82),
 \qquad LK=0,\quad HE_0=0.
\]
The program supplies each matrix and a nonsingular minor for each stated rank, with zero products checked directly. These are certificates over $\mathbb F_2$ and its extensions; they do not assert that the chosen minors are units over the integers.

For completeness, the obstruction is as follows. Border rank at most ten would supply a ten-plane $E\supset U$ with
$\operatorname{rank}\Phi_A(E)\leq4\cdot10=40$ and
$\operatorname{rank}\Phi_B(E)\leq8\cdot10=80$.
These necessary conditions follow by taking the closed proper image of the incidence of $E$ with the two determinantal rank loci; on a space spanned by ten pure elements, each wedge image has dimension at most four or eight per element. Since $K$ already has rank forty, the first condition forces every $LZ_s$ to vanish on $E$. The recorded ranks give $\ker H=\operatorname{span}E_0$, so $E$ must equal that ten-plane. Its second wedge rank is eighty-two, a contradiction. Thus $\underline{\operatorname R}(T_5)\geq11$ in characteristic two. The eleven-term finite-difference curve has integral cancellations and gives equality there.

These arguments establish lower and upper statements with their party structures and field scopes. They do not turn a three-party border degeneration into a cheaper four-party wire cover. Proposition~\ref{prop:single-layout} already excludes every one-layout Euclidean saving for a wire product, including degenerations; sums still require a quantitative selector account and an output wire identity.
